% !TeX root = ../main.tex
% !TeX encoding = UTF-8

\chapter{Paper Shield, the Post Office Horizon Scandal and the Governance Record That Destroyed Lives}

This chapter is different from the two that precede it.

Chapters 5 and 6 examined Digital Alibi failures whose consequences were institutional and financial. Aviation produced a federal mandate and an industry transformed. The 2008 financial crisis produced regulatory frameworks, institutional settlements, and a personal liability regime. The harm was systemic and economic. The individuals most directly harmed were, in the main, shareholders, mortgage holders, and economies, diffuse and statistical rather than named and specific.

The Post Office Horizon scandal is different. The individuals most directly harmed have names. Many of them went to prison. Several died before the Inquiry established that they had done nothing wrong. Sub-postmasters who had run their local branches with honesty and care for years found themselves prosecuted for theft, false accounting, and fraud on the basis of evidence that the Post Office's governance architecture, its information picture at each decision to prosecute, could not subsequently be shown to have adequately examined. The Digital Alibi gap in this case did not produce a fine or a regulatory sanction. It produced wrongful imprisonment.

This is why Chapter 6 ended where it did: the shift from governance problem to justice problem is not rhetorical. It is the actual scale difference between a case in which an organisation cannot demonstrate its decision-making was adequate and a case in which the absence of that demonstration permitted an injustice to continue for more than two decades. Boards governing AI-assisted decisions in 2026 are not making decisions about sub-postmasters. But the governance architecture failure that allowed the Post Office to continue prosecuting when it should have stopped is the same architecture failure that will allow any organisation to continue operating an inadequately governed AI system past the point at which the evidence of its inadequacy should have halted it. The Digital Alibi is not just about protecting organisations from regulatory scrutiny. At its furthest reach, it is about preventing systems from doing harm that cannot be undone.

\scenebreak

\section{The Horizon System and the Governance Position}

Horizon was an IT system developed by Fujitsu and deployed by the Post Office from 1999. It was designed to automate accounting functions in post office branches, replacing paper-based reconciliation with a digital system that tracked transactions and flagged discrepancies. The system was, by the standards of the late 1990s, a significant technology programme. Its governance, from the outset, was inadequate to the consequences its outputs would be used to produce.

The central governance claim the Post Office maintained, for more than two decades, was that Horizon was reliable. That when the system showed a financial discrepancy in a sub-postmaster's account, the discrepancy reflected a real shortfall for which the sub-postmaster was responsible. This claim was not made without basis, the Post Office had conducted internal testing, had received assurances from Fujitsu, and had governance documentation that described the system's reliability. The claim was, in governance terms, a policy position: a statement of institutional belief, supported by documentation of intent and design.

What the claim was not, and what the Horizon Inquiry has been establishing, in exhaustive public detail, since 2020, is an information picture. The Post Office's governance architecture was not designed to capture, at the moment of each decision to prosecute a sub-postmaster, the complete information picture of what was known about Horizon's reliability. The internal bug reports existed. The second opinions from technical experts questioning the system's outputs existed. The legal correspondence raising concerns about the use of Horizon evidence existed. The information was in the organisation. It was not captured in a form that established, at each specific prosecutorial decision moment, that the decision-maker had seen it, weighed it, and concluded that the prosecution should proceed despite it.

\begin{quote}
\emph{The information about Horizon's unreliability existed in the organisation. It was not captured in a form that established, at each prosecutorial decision moment, that it had been seen, weighed, and found insufficient to stop the prosecution. That is the Digital Alibi gap. And in this case, it destroyed lives.}
\end{quote}

\scenebreak

\section{What Existed, What Was Accessible, What Was Missing}

The Horizon Inquiry's methodology is, in forensic terms, the largest Digital Alibi reconstruction exercise currently under way in the United Kingdom. Sir Wyn Williams and his team are doing exactly what the Digital Alibi protocol would have made unnecessary: attempting to establish, from whatever contemporaneous evidence survived, what the Post Office's decision-makers knew at the moment of each decision to prosecute, defend, or continue opposing sub-postmasters' challenges.

What existed: the documentation of Horizon's design and intended function. The assurances from Fujitsu about system reliability. The governance committee minutes recording that Horizon's performance had been reviewed. The prosecution files compiled for each sub-postmaster case. The legal arguments advanced in court. The board papers discussing the litigation strategy. All of this existed and has been produced to the Inquiry.

What was accessible: the Inquiry has been able to examine all of this documentation. It has interviewed serving and former Post Office executives, board members, Fujitsu engineers, legal advisers, and investigators. It has reviewed thousands of documents. Everything that was retained and not destroyed has been made available.

What was missing: the information picture at each specific decision moment. When the decision was made to prosecute a specific sub-postmaster, what information about Horizon's known bugs and anomalies was before the decision-maker? When the decision was made to continue defending the Post Office's position in court against challenges to Horizon's reliability, who reviewed the technical evidence, in what form, with what authority, and on the basis of what information about the system's actual performance? When the decision was made, by a named individual, at a named moment, to proceed rather than pause, what did that individual know, and can it be shown that they knew it from a contemporaneous record rather than a reconstructed account?

The Inquiry has found, repeatedly and at every level of the Post Office's governance structure, the same answer: the information picture at each specific decision moment was not captured in a form that allows those questions to be answered from the contemporaneous record. The decisions were made. The information existed. The information picture at the decision moment does not.

\scenebreak

\section{The Horizon Scandal Through the Lens of the Digital Alibi}

Mapping the Horizon governance failure against the four Digital Alibi elements illustrates precisely how each element's absence contributed to the injustice the Inquiry is uncovering.

Decision Identity. The prosecutorial decisions, to bring charges, to continue prosecuting, to oppose appeals, to maintain the reliability of Horizon evidence in court, were not always documented with the specificity required to establish which individual made which decision on which date on the basis of which assessment of Horizon's reliability. The Post Office's prosecution function was operated by a small team whose decisions were not always individually attributed in governance records. When the Inquiry asks who decided to prosecute a specific sub-postmaster, and on what basis, the answer is frequently that the documentation shows a decision was made without establishing, with the precision the Inquiry requires, who made it, when, and with what information.

The Information Picture. The technical evidence about Horizon's known defects, the bugs reported by sub-postmasters, the Fujitsu Known Error Logs, the second opinions commissioned by legal advisers representing sub-postmasters, existed at various points in time within the Post Office's information environment. What did not exist was a governance architecture that required this information to be systematically assembled and presented at the moment of each prosecutorial decision. The decision to prosecute was taken on the basis of the Horizon discrepancy. The information about Horizon's potential to produce discrepancies through system error was not, at each prosecutorial decision moment, part of the information picture that governed that decision.

The Human Oversight Record. The Post Office's internal investigators compiled prosecution files. Legal advisers reviewed those files. Prosecution decisions were taken. But the record of human oversight, who specifically reviewed what specific technical evidence about Horizon's reliability at what stage of each prosecution decision, is, in many cases, irrecoverable. The oversight existed. The record of its content and quality at each specific decision moment did not.

The Accountability Chain. The Inquiry has established that significant decisions about prosecutorial strategy, about how to respond to challenges to Horizon's reliability, and about the disclosure of technical evidence were made at varying levels of the Post Office's executive and governance structure. The accountability chain for those decisions, which named individual accepted responsibility for which specific decision, documented at the moment the decision was made, is, in multiple cases, unestablished by the contemporaneous record.

\begin{longtable}[]{@{}
  >{\raggedright\arraybackslash}p{(\columnwidth - 0\tabcolsep) * \real{1.0000}}@{}}
\toprule\noalign{}
\begin{minipage}[b]{\linewidth}\raggedright
\textbf{THE HORIZON DIGITAL ALIBI FAILURE, ELEMENT BY ELEMENT}
\end{minipage} \\
\midrule\noalign{}
\endhead
\bottomrule\noalign{}
\endlastfoot
Decision Identity: prosecutorial decisions not individually attributed in contemporaneous governance records. The Inquiry asks ``who decided?'', the documentation answers ``a decision was made.'' \\
The Information Picture: technical evidence of Horizon's defects existed in the organisation but was not systematically assembled at each prosecutorial decision moment. The decision was taken on the discrepancy; the information about the system's capacity to generate spurious discrepancies was not part of the contemporaneous decision record. \\
The Human Oversight Record: oversight existed but its content, what was reviewed, by whom, at what level of technical scrutiny, is irrecoverable from the contemporaneous documentation for many individual cases. \\
The Accountability Chain: responsibility for decisions about prosecution strategy, disclosure, and response to challenges distributed across the executive and governance structure without contemporaneous individual attribution sufficient to sustain accountability. \\
\end{longtable}

\scenebreak

\section{Policy Position vs.~Information Picture: The Core of the Governance Failure}

The Post Office's governance position on Horizon's reliability was a policy position. It was a statement, maintained consistently and communicated with institutional authority, that the system was reliable and that discrepancies reflected genuine sub-postmaster accounting failures. That policy position was documented, reviewed by governance bodies, and supported by legal advice. It had all the formal attributes of a governance decision.

What it did not have was an information picture at the level of each individual prosecutorial decision. The policy said Horizon was reliable. The information picture at each prosecution decision would have required establishing, for the specific sub-postmaster, on the specific date, whether the specific discrepancy had been assessed against the specific known bugs and anomalies in the system at that point in time, and whether a named individual had concluded, on the basis of that specific assessment, that the discrepancy reflected genuine fraud rather than system error.

The policy existed. The information picture did not. And because the information picture did not exist, the policy could be maintained even as the evidence of its inaccuracy accumulated within the organisation. A governance architecture designed to capture the information picture at each prosecutorial decision moment would have made the maintenance of an inaccurate policy increasingly difficult to sustain, because each decision moment would have required a contemporaneous record of what was known about Horizon's reliability, and the accumulation of that record would have made the discrepancy between policy and reality visible in the governance record itself.

\begin{quote}
\emph{A governance architecture that captures the information picture at each decision moment makes an inaccurate policy increasingly difficult to sustain. The Horizon governance architecture did not capture that picture. The policy was sustained for more than two decades while the evidence of its inaccuracy accumulated alongside it.}
\end{quote}

This is the most important governance lesson the Horizon scandal offers for AI governance in 2026. Boards governing AI-assisted decisions will, in many cases, maintain policy positions about their systems' reliability, fairness, and governance adequacy. Those positions will be documented, reviewed, and supported by governance frameworks. They will have all the formal attributes of a governance decision. Without a Digital Alibi architecture that captures the information picture at each output moment, those policy positions can be maintained even as the evidence of their inaccuracy accumulates. The governance record will show the policy. It will not show the information picture at each decision moment. And the gap between those two things is where the harm occurs.

\scenebreak

\section{The Inquiry as Digital Alibi Reconstruction}

The Horizon Inquiry is, in governance terms, an exercise in retroactive Digital Alibi reconstruction. Sir Wyn Williams and his team are attempting to establish, from board papers, internal emails, Fujitsu correspondence, prosecution files, legal advice, and the testimony of witnesses, what the Post Office's decision-makers knew at each specific decision moment, and whether what they knew was sufficient to justify the decisions they made.

The reconstruction is difficult, expensive, and incomplete. It has required more than four years of public hearings and the examination of hundreds of thousands of documents. And it will remain incomplete, because the contemporaneous documentation for many of the most consequential decision moments simply does not exist. The Inquiry can establish patterns. It can establish what was in the organisation's information environment at broad points in time. What it cannot establish, for many specific decisions, is the information picture at the exact moment that decision was made. The contemporaneous record that would have made that reconstruction unnecessary was never created. The reconstruction is the cost of its absence.

The sub-postmasters who were prosecuted on the basis of a governance architecture that did not capture the information picture at each decision moment have spent years of their lives under criminal conviction, in some cases in prison, in the certainty that the institution prosecuting them knew more than it was acknowledging. Several died before the Inquiry established their innocence. Their families carry the consequences of a governance failure that was not malicious in its origins but became, through the absence of the information picture that would have forced a reckoning, an injustice sustained by institutional momentum and the absence of contemporaneous evidence sufficient to challenge it.

\begin{quote}
\emph{The Horizon Inquiry is doing in public what the Digital Alibi protocol prevents in private: retroactively reconstructing what an organisation's decision-makers knew at each specific moment of a consequential decision. The reconstruction is incomplete, expensive, and too late for those who needed it most.}
\end{quote}

\scenebreak

\section{The AI Governance Parallel, Before the Inquiry Arrives}

The Horizon parallel for AI governance is not hypothetical. It is structural. Organisations deploying AI in material decision processes, in credit, underwriting, hiring, procurement, benefits assessment, criminal risk scoring, are, in the absence of a Digital Alibi architecture, creating the conditions that produced the Horizon injustice in a different technological domain.

The scale is different. The technology is different. The decisions are different. But the governance architecture failure is the same: decisions made on the basis of AI system outputs, with a policy position asserting the system's reliability and adequacy, without a contemporaneous information picture at each output moment that would allow a subsequent examination to establish what the system knew, what it was not given, who reviewed its output, and who was accountable at the decision moment.

The sub-postmasters whose accounts showed Horizon discrepancies could not, in most cases, challenge the reliability of the evidence against them because the Post Office's governance architecture did not capture the information that would have made that challenge possible. The individuals whose credit applications, benefits assessments, or employment screenings produce AI-assisted adverse decisions face the same structural challenge: the information picture at the moment the system produced its output is not captured in a form that would allow the decision to be substantively examined. The policy says the system is reliable. The information picture at the output moment does not exist. The challenge cannot be mounted.

The Horizon Inquiry is making this dynamic visible in the most public and consequential way possible. Every NED who has watched its proceedings, and in the United Kingdom, it is difficult to have not, has seen, in real time, what the absence of an information picture at each decision moment produces when examined by a public inquiry with the power to compel evidence and hear testimony. They have seen the institutional destruction, the personal accountability, and the human cost.

The question the Horizon Inquiry poses for every board governing AI-assisted decisions is not ``could this happen to us?'' It is: ``are we building the governance architecture that would make it impossible for it to happen to us?'' The answer, for most organisations, is no. The architecture exists. The Digital Alibi protocol can be built. Chapter 7 exists to make the cost of not building it impossible to ignore.

\scenebreak

\section{The Boardroom Mirror: Where This Happens Next}

The Post Office Horizon scandal is not a story about a unique organisational failure or an unusually corrupt institution. It is a story about what happens when a governance architecture is built to maintain a policy position rather than capture an information picture. That architecture is not unusual. It is the default architecture of most organisations deploying technology in consequential decision processes. The question is not whether your organisation's AI governance resembles the Post Office's Horizon governance. The question is which element of the parallel you are most willing to examine.

Consider a healthcare trust that deploys an AI clinical decision support system to assist in triaging patient referrals for elective procedures. The system has been validated, its performance metrics reviewed by the clinical governance committee, and its deployment approved by the board with a policy position that the system supports, not replaces, clinical judgement. That policy position is documented, periodically reviewed, and consistent with the trust's stated commitment to responsible AI use.

\begin{longtable}[]{@{}
  >{\raggedright\arraybackslash}p{(\columnwidth - 0\tabcolsep) * \real{1.0000}}@{}}
\toprule\noalign{}
\begin{minipage}[b]{\linewidth}\raggedright
\textbf{THE AI ``HORIZON'' SCENARIO, HEALTHCARE TRIAGE}
\end{minipage} \\
\midrule\noalign{}
\endhead
\bottomrule\noalign{}
\endlastfoot
The policy position: ``Our AI supports clinical judgement in triage. It is validated and governed.'' \\
The information picture at each output moment: the system was trained on historical referral data from a predominantly urban patient population. It was not given socioeconomic deprivation indices at the point of each triage assessment. Its outputs for patients from specific postcodes systematically underweight clinical urgency, not because the algorithm is biased in design, but because the data it was not given at the output moment is predictively significant. \\
What the governance record shows: policy compliance, validation reports, committee minutes, performance dashboards. \\
What the governance record does not show: for each specific triage decision, what data the system was not given, whether anyone reviewed that absence, and who was named and accountable for the decision to deploy a system with that specific negative space. \\
The accumulation: adverse triage decisions concentrate in specific patient populations. No single decision triggers a governance review. The pattern is visible only in the aggregate, and only if someone looks for it with the information picture at the output moment available. \\
The Horizon parallel: the Post Office's governance record showed policy compliance. It did not show what was known about Horizon's defects at each prosecutorial decision moment. The trust's governance record shows policy compliance. It does not show what was not given to the system at each triage decision moment. The governance architecture is the same. The harm accumulates the same way. \\
\end{longtable}

The difference between this scenario and Horizon is not structural. It is temporal. The Post Office's governance failure accumulated over two decades before a public inquiry began to reconstruct what the information picture at each decision moment would have shown. The trust's governance failure is accumulating now. The Digital Alibi protocol, output-moment capture of what the system knew, what it was not given, who reviewed its output, and who was accountable, is the architecture that stops the accumulation. It does not prevent the policy from being wrong. It makes the wrong policy impossible to sustain in the absence of the information picture that would show it.

\scenebreak

Chapter 7 has examined the Digital Alibi gap in its most consequential human form, an injustice sustained by the absence of an information picture at each decision moment, uncovered by a public inquiry that is still reconstructing what should have been captured as events occurred.

Chapter 8 turns to the Space Shuttle Challenger disaster of 1986. It is a different kind of case: not a sustained institutional injustice over decades, but a single catastrophic decision made in real time under pressure, with adequate information present in the system but absent from the information picture at the moment of the decision. The Challenger disaster produced the most studied decision failure in engineering history, and its forensic anatomy reveals something that the Horizon case does not: what happens when the information picture is deliberately structured to exclude the information that should have stopped the decision.
